\noindent{\bf Notation.}~We use the following notation in the paper. Given a matrix $\mA$, let $\emA_{i,j}$, $\mA_{i,:}$, and $\mA_{:, j}$ denote its $(i,j)$-th entry, $i$-th row, and $j$-th column, respectively. We use $\norm{\mA}_p$ to denote the entry-wise $\ell^p$ norm of $\mA$. Let $\sfmx[\cdot]$ be the softmax operator, which takes a matrix as input and applies softmax operation to each column of the matrix, which results in a column stochastic matrix, i.e., a matrix that has non-negative entries with each column summing to 1. We similarly define $\hdmx[\cdot]$ to be the hardmax operator, which outputs the one-hot representation of the $\argmax$ entry for each column of the input matrix. If there are $k$ $\argmax$ entries, then the output is $1/k$ for such entries.
We use $\vone_n$ to denote a vector of length $n$ whose entries are all $1$. We denote the 0-1 indicator function by $\indic{\cdot}$. We use $d$ and $n$ to denote the embedding dimension and the sequence length, respectively. We assume throughout that $n \geq 2$, as the Transformers reduce to residual networks when $n=1$.

\vspace*{-5pt}
\section{Transformer networks}\label{sec:transformerblocks}
\vspace*{-5pt}

%\subsection{Notation}
%Before we introduce the necessary background material on Transformer models, a short note on the notation used throughout this paper is in order. For $a, b \in \naturals$, where $a \leq b$, let $[a] \defeq \{1, 2, \dots, a\}$ and $[a:b] \defeq \{a, a+1, \dots, b\}$. We use $\vone_n$ to denote a vector of length $n$ whose entries are all $1$. We denote the 0-1 indictator function by $\indic{\cdot}$. For $t \in \reals$, $\relu(t) = \max \{t, 0\}$ denotes the rectified linear unit. The notation $\relu(\cdot)$ naturally extends to the multi-dimensional arguments by applying the operation in an entry-wise manner. 

%Similarly, we define the {\em hardmax} operator $\hdmx[\cdot]$, which transforms each column of a matrix to the one-hot representation of the argmax entry in that column. In the event when an input column has multiple (say $k$) argmax entries, the output has value $1/k$ at each of these entries.

% Transformer networks -- the main focus of this paper -- are essentially sequence-to-sequence functions mapping 
% $\reals^{d \times n}$ to $\reals^{d \times n}$.
%The main subject of study in this paper is Transformer networks as sequence-to-sequence functions mapping $\reals^{d \times n}$ to $\reals^{d \times n}$.
%Here, $d$ and $n$ denote the embedding dimension and the sequence length, respectively, which are both fixed quantities for a network. For natural language processing applications, we can view the input $\mX \in \reals^{d \times n}$ as a sequence of $n$ tokens, with each token embedded in a $d$-dimensional space. %We call a function $f : \reals^{d \times n} \to \reals^{d \times n}$ to be {\em permutation equivariant} if for any permutation matrix $\mP$, we have $f(\mX \mP) = f(\mX) \mP$; i.e., if we permute the columns of $\mX$, then the columns of $f(\mX)$ are permuted in the same way.

%\subsection{Transformer block and Self Attention}
%\label{sec:transformerblocks}


A Transformer block is a sequence-to-sequence function mapping $\reals^{d \times n}$ to $\reals^{d \times n}$. It consists of two layers: a self-attention layer and a token-wise feed-forward layer, with both layers having a skip connection. More concretely, for an input $\mX \in \reals^{d \times n}$ consisting of $d$-dimensional embeddings of $n$ tokens, a Transformer block with {\em multiplicative} or {\em dot-product} attention~\citep{luong2015multiplicativ} consists of the following two layers\footnote{In our proof we use bias vectors $\vb_Q^i$ for query projections in attention layers. We omit them here for brevity.}:
\begin{align}
    \label{eq:attn}
    \attn(\mX) &= \mX + \sum\nolimits_{i=1}^h \mW_O^i \mW_V^i \mX \cdot \sfmx [(\mW_K^i \mX)^T \mW_Q^i \mX],\\
    \label{eq:ff}
    \ff(\mX) &= \attn(\mX) + \mW_2 \cdot \relu (\mW_1 \cdot \attn(\mX) + \vb_1 \vone_n^T) + \vb_2 \vone_n^T,
\end{align}
where $\mW_O^i \in \reals^{d \times m}$, $\mW_V^i, \mW_K^i, \mW_Q^i \in \reals^{m \times d}$, $\mW_2 \in \reals^{d \times r}, \mW_1 \in \reals^{r \times d}, \vb_2 \in \reals^{d}, \vb_1 \in \reals^{r}$, and $\ff(\mX)$ is the output of the Transformer block. The number of heads $h$ and the head size $m$ are two main parameters of the attention layer; and $r$ denotes the hidden layer size of the feed-forward layer.

Here, we would like to point out that our definition of the self-attention layer \eqref{eq:attn} is an equivalent reformulation of \citep{vaswani2017attention}, where they concatenate attention heads and multiply a matrix $\mW_O \in \reals^{d \times mh}$ to the concatenation.
One difference in our setup is the absence of layer normalization, which simplies our analysis while preserving the basic architecture of the Transformer.

%\subsection{Transformer networks and their permutation equivariance}
%\label{sec:permequiv}

We define the Transformer networks as the composition of Transformer blocks. The family of the sequence-to-sequence functions corresponding to the Transformers can be defined as:
\begin{align*}
%\label{eq:transformer-fclass}
    \mc T^{h,m,r} 
    &\defeq 
    \{ g: \reals^{d \times n} \to \reals^{d \times n} \mid 
    \text{$g$ is a composition of Transformer blocks $t^{h,m,r}$'s} \}.
\end{align*}
where $t^{h,m,r}:\reals^{d \times n} \to \reals^{d \times n}$ denotes a Transformer block defined by an attention layer with $h$ heads of size $m$ each, and a feed-forward layer with $r$ hidden nodes. 

We say that a function $f : \reals^{d \times n} \to \reals^{d \times n}$ is {\em permutation equivariant} if for any permutation matrix $\mP$, we have $f(\mX \mP) = f(\mX) \mP$; i.e., if we permute the columns of $\mX$, then the columns of $f(\mX)$ are permuted in the same way. A Transformer block is permutation equivariant, which we formally prove in Section~\ref{sec:proof_equivariance}. This consequently establishes the permutation equivariance of the class $\mc T^{h,m,r}$. 
\begin{claim}
\label{claim:equivariance}
A Transformer block $t^{h, m, r}$ defines a permutation equivariant map from $\reals^{d \times n}$ to $\reals^{d \times n}$.
\end{claim}

As seen in above, both layers (cf.~\eqref{eq:attn} and \eqref{eq:ff}) of a Transformer block employ parameter reuse/sharing, because each token/column undergoes the same transformations (e.g., $\mW_Q^i$, $\mW_K^i$, or $\mW_1$) regardless of its position.
Moreover, interactions between tokens can only be captured through pairwise dot-products in the softmax operator $\sfmx[\cdot]$~(cf.~\eqref{eq:attn}).
Given such limitations in a single Transformer block's representation power, it is not obvious what kinds of sequence-to-sequence functions $\mc T^{h,m,r}$ can approximate; we provide the answer to this question in the next section.

%In order to endow the Transformer networks with the ability to capture the information about the position of tokens in the input sequence, it is a common practice to add positional embeddings $\mE \in \reals^{d \times n}$ to the input sequence before feeding it to the Transformer network~ \citep{vaswani2017attention, devlin2018bert}. The Transformed networks combined with the positional embeddings significantly broadens the set of sequence-to-sequence functions defined by such networks as it removes the restriction of permutation equivariant. This motivates us to consider the following extension of~\eqref{eq:transformer-fclass}: 


%It is natural to ask how rich the function classes $\mc T^{h, m, r}$ and $\mc T^{h, m, r}_{\rm P}$ are, i.e., {\em what are all the sequence-to-sequence functions that we can hope to approximate by Transformer networks?}